\section{Integrability of Gauduchon determinant measures}
\label{sec:gaud-determinant-integrability}

The lower bound for equation normalizations gives a uniform $L^q$ estimate for the determinant measures of all normalized Gauduchon currents. Every exponent below $n/(n-1)$ is available. At the endpoint, the optimal constant is exactly $v_-(\omega)^{-1/n}$, linking integrability of the dual tests to the lower mass of the background.

Write $\Gamma_0(T)$ for the determinant measure in \eqref{eq:weak-gaud-functional} with $F=0$.

\begin{theorem}\label{thm:gaud-determinant-subcritical}
For every $1<q<n/(n-1)$ there is $C_q>0$ such that every $T\in\mathscr G_\omega$ satisfies
\[
 \Gamma_0(T)=k_T\,\mathrm{d}V,\qquad k_T\geq0,\qquad
 \|k_T\|_{L^q(\mathrm{d}V)}\leq C_q.
\]
No bounded mass or positive lower-volume hypothesis is required.
\end{theorem}
\begin{proof}
Fix $p>1$ and take any smooth strictly positive function $h$. The smooth equation gives $\Omega_h=\omega+\ddc u_h>0$ and $c_h>0$ such that
\[
 \Omega_h^n=c_hh^n\,\mathrm{d}V.
\]
By \cref{thm:lp-potential-estimate}, the density of this bounded solution satisfies
\[
 c_h\|h^n\|_{L^p(\mathrm{d}V)}\geq c_p>0.
\]
Pointwise matrix arithmetic--geometric mean, applied to $\Omega_h$ and the coefficient density of $T$, gives the measure inequality
\[
 c_h^{1/n}h\,\Gamma_0(T)\leq\Omega_h\wedge T.
\]
Integrating and using $\ddc T=0$ yields
\[
 \int_Xh\,\mathrm{d}\Gamma_0(T)
 \leq c_h^{-1/n}\int_X\Omega_h\wedge T
 =c_h^{-1/n}
 \leq c_p^{-1/n}\|h\|_{L^{np}(\mathrm{d}V)}.
\]
Approximation by smooth strictly positive functions extends this bound to every continuous nonnegative $h$. Applied to $|h|$, it bounds the positive measure functional on all continuous functions by the $L^{np}$ norm. Density of continuous functions in $L^{np}$ and the usual duality of these spaces show that the measure has a density in $L^{np/(np-1)}$, with norm at most $c_p^{-1/n}$. As $p$ ranges over $(1,+\infty)$, the exponent $np/(np-1)$ ranges over $(1,n/(n-1))$. This proves the result.
\end{proof}

In particular, the determinant of the matrix density vanishes on the singular part of the trace measure: if $\sigma_T=\sigma_{T,\mathrm{ac}}+\sigma_{T,\mathrm{s}}$ is its Lebesgue decomposition, then
\[
 \det P_T=0\quad\sigma_{T,\mathrm{s}}\text{-almost everywhere}.
\]
Indeed $f_0=(\mathrm{d}V/\omega^n)^{1/n}$ is everywhere positive and the determinant-weighted measure is absolutely continuous. This assertion concerns the determinant; it does not remove a possible singular part of the current itself.

\Needspace{12\baselineskip}
\begin{theorem}\label{thm:gaud-determinant-endpoint}\index{Gauduchon metric!determinant integrability}
With the extended value $+\infty$ allowed, one has
\[
 \sup_{T\in\mathscr G_\omega}
       \|k_T\|_{L^{n/(n-1)}(\mathrm{d}V)}=v_-(\omega)^{-1/n},
\]
where the right side is interpreted as $+\infty$ when $v_-(\omega)=0$. The same supremum is obtained by restricting to smooth strictly positive normalized Gauduchon forms. Thus a uniform estimate at the endpoint exponent is equivalent to $v_-(\omega)>0$.
\end{theorem}
\begin{proof}
Suppose first that $v_-(\omega)>0$. In the smooth equation of the preceding proof, the definition of $v_-$ gives
\[
 c_h\int_Xh^n\,\mathrm{d}V=\int_X\Omega_h^n\geq v_-(\omega).
\]
The same pointwise inequality and Stokes therefore give
\[
 \int_Xh\,\mathrm{d}\Gamma_0(T)
      \leq v_-(\omega)^{-1/n}\|h\|_{L^n(\mathrm{d}V)}.
\]
Duality proves the asserted upper bound for every current in $\mathscr G_\omega$.

\Needspace{8\baselineskip}
For the reverse bound, let $g$ be any smooth positive probability density relative to $\mathrm{d}V$, and solve
\[
 \Omega_g^n=c(g)g\,\mathrm{d}V.
\]
Let $T_g$ be the trace-normalized $(n-1)$st power of the Gauduchon metric conformal to $\Omega_g$. Equality in the pointwise matrix inequality gives
\[
 \int_Xg^{1/n}\,\mathrm{d}\Gamma_0(T_g)=c(g)^{-1/n}.
\]
Since $\|g^{1/n}\|_{L^n}=1$, its endpoint dual norm is at least $c(g)^{-1/n}$. By \cref{thm:gaud-volume-mass-extrema},
\[
 \inf_{g>0\ \mathrm{smooth},\ \int_Xg\,\mathrm{d}V=1}c(g)=v_-(\omega).
\]
Taking the supremum over these smooth $T_g$ proves the matching lower bound. If $v_-(\omega)=0$, the same argument makes this supremum infinite. This also proves that allowing nonsmooth currents does not change it.
\end{proof}

The endpoint norm is independent of the smooth probability measure used to express it, as the formula shows. Subcritical constants may depend on that choice. Neither assertion identifies $v_-$ with the upper mass $v_+$.

\begin{corollary}\label{cor:gaud-adjoint-lq}
Let $\mathcal A_0(F)$ be the supporting probabilities in \cref{thm:gaud-current-all-adjoints}. For every $R>0$ and $1<q<n/(n-1)$ there is $C_{R,q}$ with the following property. If $F\in C(X)$ and $\|F\|_\infty\leq R$, every $\nu\in\mathcal A_0(F)$ has a density $a$ satisfying
\[
 \nu=a\,\mathrm{d}V,\qquad\|a\|_{L^q(\mathrm{d}V)}\leq C_{R,q}.
\]
If $v_-(\omega)>0$, this holds also for $q=n/(n-1)$.
\end{corollary}
\begin{proof}
Write $M(F)=\exp(-B_0(F)/n)$ and represent $\nu$ by a maximizing current $T\in\mathscr M(F)$, using \cref{thm:gaud-current-all-adjoints}. Then
\[
 \nu=\frac{\mathrm{e}^{F/n}k_T}{M(F)}\,\mathrm{d}V,
 \qquad M(F)\geq \mathrm{e}^{-R/n}M(0)>0.
\]
The two preceding theorems give the bound, for example with constant
$\mathrm{e}^{2R/n}M(0)^{-1}C_q$ in the subcritical case. Use the endpoint constant when it is finite.
\end{proof}

The endpoint fails for the full normalized Gauduchon family on the threefold of \cref{thm:vanishing-lower-mass}, while every subcritical estimate above remains valid. For a fixed smooth right-hand side, \cref{thm:weak-gaud-compact-duality} gives a unique smooth maximizing current, whose determinant density is bounded. This does not give uniform higher moments along sequences approaching its value; \cref{prop:gaud-determinant-concentration} provides a counterexample.

\Needspace{13\baselineskip}
The endpoint can be available on a manifold with infinite upper mass. On the Hopf threefold of \cref{chap:hopf}, take its Gauduchon metric $\omega$, for which $v_-(\omega)=6$ and $v_+(\omega)=+\infty$, and set $\mathrm{d}V=\omega^3/6$. The normalized Gauduchon form $T_0=\omega^2/6$ has $\sigma_{T_0}=\mathrm{d}V$ and $P_{T_0}=I$. Hence $k_{T_0}=6^{-1/3}$, and
\[
 \sup_{T\in\mathscr G_\omega}\|k_T\|_{L^{3/2}(\mathrm{d}V)}=6^{-1/3}.
\]
This gives an attained endpoint determinant estimate alongside failure of bounded mass. The estimate controls these finite Gauduchon measures, not the total mass of arbitrary nonlinear products.
